\section{Width and coefficient size}\label{sec:width-coefficient}

The preceding encoder has two independent design parameters: the number
of preceding states used by a scan, and the coefficient alphabet.
Increasing either parameter preserves the vector lengths and the rate.
Their costs differ. A wider scan performs more arithmetic and reads more
coefficients. A larger alphabet can retain the same machine instruction
sequence until a coefficient or metadata representation changes.

\subsection{A common family}

Fix $p=2^{128}-159$, $k=2^{20}$, and $n=2k$. Unless stated otherwise,
the local source map in this section is the extended Newton code
\begin{equation}\label{eq:extended-newton-eight}
 f_x(t)=\sum_{j=0}^{7}\binom tj x_j,\qquad
 B_8(x)=(f_x(0),\ldots,f_x(14),x_7).
\end{equation}
Apply $B_8$ independently to every eight input symbols. This gives a
length-$n$ vector. The local code is $[16,8,9]$ MDS: if $x_7\ne0$,
at most seven finite evaluations vanish; if $x_7=0$, the polynomial
has degree at most six. In either case a nonzero block has weight at
least nine. Its first eight evaluations also establish injectivity.
Finite differences evaluate the map with 77 additions per block,
including only updates that can reach a remaining output.

For $1\le w<n$, a width-$w$ scan $A_{w,a}$ maps $v$ to $y$ by
\begin{equation}\label{eq:width-scan}
 y_i=v_i+\sum_{j=1}^{\min(w,i)}a_{i,j}y_{i-j},
 \qquad 0\le i<n.
\end{equation}
Here $1\le q<p$, and every coefficient $a_{i,j}$ is sampled independently
from $\{1,\ldots,q\}$. For $L$ rounds, sample independent uniform
permutation matrices $P_1,\ldots,P_L$, independently of all coefficients.
With column-vector notation, the encoder is
\begin{equation}\label{eq:width-family}
 E=A_{w,a_L}P_L\cdots A_{w,a_1}P_1 B_8.
\end{equation}
Every scan is invertible. The encoder therefore has dimension $k$ for
every setup, and every vector after the source has length $n$.
The output is the complete final vector.

We also consider an additional scan $A_{t,b}E$, without another
permutation. We call $t$ the final unpermuted width; $t=0$ denotes its
absence. This parameter is stated separately from $L$ and $w$.
Each distance guarantee below bounds
\[
 \Pr_{P,a,b}\bigl[\exists x\in\F_p^k\setminus\{0\}:
                         \wt(E x)<d_0\bigr],
\]
with $E$ replaced by $A_{t,b}E$ when applicable. Setup precedes the
choice of message. The prime is fixed before setup.

\subsection{Two width-one scans with small coefficients}

The all-one case already has a complete baseline with the smaller local
$[8,4,5]$ source from Section~\ref{sec:small-scans}. Set every scan
coefficient to one and sample only the two independent uniform
permutations. At the same $k$ and $n$, the code has minimum distance at
least 75498 ($>3.6\%$), except with probability less than $2^{-31.61}$.
This bound holds for each fixed prime $p>6$, including the target prime.
It covers every nonzero message after setup. Encoding uses exactly
$3.875n-2=8126462$ field additions. A support-only count proves this
baseline; it uses neither random multipliers nor a union over field values.
The complete rational certificate is included with the companion results.
This historical bound does not reach our comparison target of 40 bits
of failure margin. With the larger local block and refined analysis,
the $q=1$ row below now gives distance above $10.7\%$ and failure below
$2^{-46.39}$. This supplies a matched all-one baseline for the other
coefficient alphabets.

A different local block gives a complementary point. The eight-input
subset-sum construction of Appendix~\ref{app:subset-sum-gv} reaches
distance above $11.4\%$ with failure below $2^{-24.26}$ and only
$3.3125n-2$ additions. Both scans still have coefficient one. This
point uses a different source catalog and a smaller failure margin
than the MDS-source comparison below.

The larger local block and a sharper treatment of dependent equations
give the following finite guarantees. No final unpermuted scan is used.

\begin{theorem}\label{thm:width-one-small}
For each row of Table~\ref{tab:width-one-small}, the encoder
in~\eqref{eq:width-family} with $L=2$ and $w=1$ has minimum distance
at least the listed integer, except with probability less than the
listed setup failure bound.
\end{theorem}

\begin{table}[t]
\centering
\caption{Complete width-one guarantees for the local $[16,8,9]$ source,
two permutations, and no final unpermuted scan. Failure exponents are
rounded down. These are achieved points, not optimal frontiers.}
\label{tab:width-one-small}
% Generated by build_result_inventory.py; do not edit by hand.
\begin{tabular}{rrrr}
\toprule
$q$ & Minimum distance & Relative bound & Failure below\\
\midrule
1 & 224396 & $>10.7\%$ & $2^{-46.39}$ \\
2 & 209716 & $>10\%$ & $2^{-43.29}$ \\
4 & 335545 & $>16\%$ & $2^{-44.99}$ \\
8 & 461374 & $>22\%$ & $2^{-52.31}$ \\
14 & 534774 & $>25.5\%$ & $2^{-53.54}$ \\
16 & 545260 & $>26\%$ & $2^{-54.31}$ \\
32 & 608175 & $>29\%$ & $2^{-56.32}$ \\
64 & 629146 & $>30\%$ & $2^{-57.88}$ \\
128 & 639632 & $>30.5\%$ & $2^{-58.86}$ \\
\bottomrule
\end{tabular}

\end{table}

We describe the improved bound underlying the finite evaluation.
Fix a source support and let $D$ be the dimension of its shortened
message space. As in Section~\ref{sec:small-scans}, expose tap
coefficients while retaining a single space $W$ of imposed zero
equations across both scans. For $\tau\ge1$, attach the potential
$\tau^{-\rank W}$ to a surviving branch.

There are two feedback states: zero and nonzero. Let $x$ mark nonzero
inputs, $y$ mark nonzero outputs, and $v_0,v_1>0$ be continuation bounds.
The following positive homogeneous operator bounds the two states:
\begin{align}
 T_0(v)&=v_0+xyv_1,\nonumber\\
 T_1(v)&=yv_1+x\max\left\{
        yv_1+\frac{v_0}{\tau},\;
        \left(1-\frac1q\right)yv_1+\frac{v_0}{q}
                         \right\}.\label{eq:width-one-quotient}
\end{align}
To justify the maximum, consider a nonzero input form and a nonzero
feedback form in the quotient by $W$, before exposing the fresh tap.
If these forms are independent, requiring a zero output adds one
independent equation. This gives the first alternative. If the forms
are dependent, at most one coefficient produces the zero form.
The zero and nonzero outcomes are then mutually exclusive, giving the
second alternative when this coefficient belongs to the alphabet.
If it does not, the first alternative bounds the nonzero outcome.
The earlier bound charged a return without retaining this exclusivity.
For $q=1$, the second alternative is simply $v_0$: a dependent zero
output is deterministic. The argument still applies because that zero
output and a nonzero output cannot occur in the same realization.
Only the two permutations are random in this case.

Normalize $v_0=1$, set $\lambda=\max_i T_i(v)/v_i$, and
$C=\max_i1/v_i$. Monotonicity and homogeneity imply
$[T^n(\mathbf1)]_0\le C\lambda^n$. For sparse supports we instead use
\[
 [T^n(\mathbf1)]_0\le
 \lambda^{n-16}\max_i\frac{[T^{16}(\mathbf1)]_i}{v_i}.
\]
The terminal iterations are evaluated exactly over the rationals.

For this source, put
\[
 Q_8(z,\tau)=1+\sum_{d=9}^{16}\binom{16}{d}z^d\tau^{d-8},
 \qquad F_s(\tau)=[z^s]Q_8(z,\tau)^{n/16}.
\]
Fix source weight $s$, middle weight $h$, and forbidden final weight
$e$. Choose a positive witness for each scan, with $0<y_2\le1$ and
the same $\tau$ for both scans. Permutation averaging gives the bound
\begin{equation}\label{eq:width-one-profile}
 \frac{F_s(\tau)}{\tau}
 \frac{C_1\lambda_1^n C_2\lambda_2^n}
 {\binom ns\binom nh\,x_1^s(y_1x_2)^h y_2^e}.
\end{equation}
The divisor $\tau$ accounts for the one remaining dimension of a
nonzero message. Each permutation contributes the reciprocal number
of supports of its incoming weight. No union over field-valued messages
is used.

The finite calculation sums over all $h\le\min(n,2e)$ and
$s\le\min(n,2h)$. It treats the possible source weights below 32
using exact source coefficients and the terminal bound for $q\ge2$.
For $q=1$, an exact support-moment bound is stronger in this range;
Appendix~\ref{app:all-one} gives its finite evaluation. The remaining
weights use $F_s\le Q_8(z,\tau)^{n/16}/z^s$ and convex bounds on
regions of the $(s,h)$ domain. Every region includes its full number
of integer profiles. Positive witnesses are proposed numerically and
replayed with rational arithmetic and directed logarithm bounds.
The replay checks coverage and sums upward-rounded contributions.
Table~\ref{tab:width-one-small} records the resulting complete sums.
The detailed evaluation and its saved witnesses accompany the paper;
the two-state recurrence above supplies the rows with $q\ge2$.
For $q=1$, it proves above $6.6\%$. Jointly exposing the first two
active inputs of each run, and counting proportional pairs that cannot
cancel, improves this to $10\%$:
Appendix~\ref{app:fresh-pairs} gives the three-state macro bound and
its finite sampling correction. Counting a return together with the next
complete pair improves this to the displayed $10.7\%$;
Appendix~\ref{app:joint-renewal} gives the joint continuation bound and
its uniform rational verification.

All rows have $2n-2$ tap contributions and $6.8125n-2$ field additions,
including the source, before coefficient multiplication. For $q=1$,
the latter is the complete cost: exactly 14286846 field additions.
For $q\ge2$, a native implementation may use $2n-2$ word products.
With coefficients in $\{1,2\}$, each product requires
at most one doubling. The entire encoder then uses at most
$8.8125n-4$ additions, with expectation $7.8125n-3$ over setup.
For general $q$, the binary-method expectation is
\[
 4.8125n+2(n-1)\left(1+\frac1q\sum_{a=1}^q
 [\lfloor\log_2a\rfloor+\operatorname{popcount}(a)-1]\right).
\]
Two length-$n$ field buffers occupy 64 MiB. A 32-bit word stores a
21-bit permutation index and each listed coefficient-minus-one, giving
16 MiB for the two metadata arrays. This 80 MiB model excludes separately
retained input and validation storage. The timings in
Table~\ref{tab:small-timing} use a smaller local source and do not
measure the constructions in Table~\ref{tab:width-one-small}.

\subsection{Wider scans and a third permutation}

Larger widths also admit complete finite-coefficient bounds. The current
high-distance points use an additional unpermuted scan:
\begin{center}
\begin{tabular}{rrrrrr}
\toprule
$L$ & $w$ & $q$ & $t$ & Minimum distance & Failure below\\
\midrule
2 & 10 & $2^{63}$ & 1 & 964690 ($>46\%$) & $2^{-40.2205}$\\
3 & 10 & $2^{63}$ & 2 & 985662 ($>47\%$) & $2^{-40.0476}$\\
\bottomrule
\end{tabular}
\end{center}
The underlying argument permits any fixed local $[16,8,9]$ MDS map,
including~\eqref{eq:extended-newton-eight}. The original arithmetic
artifacts used sixteen finite evaluations; substituting the extended
map changes the source cost from 105 to 77 additions per block.
These points use $21n-111$ and $32n-168$ weighted tap contributions,
respectively. They have not been benchmarked in these configurations.

The proof first treats each tap before the final scan as an independent
indeterminate. Outside a bad event over permutations, a combinatorial
rank argument proves that every sufficiently small set of parity-check
columns is independent. Specializing a group of taps from one column
of an inverse scan is an affine rank-one update. It loses at most one
rank, with conditional probability at most $1/q$.
The final scan contributes zero equations on the shortened entering
space. An active zero either reduces this space's dimension or costs
another conditional factor $1/q$. Counting the earlier rank losses and
these redundant zero equations together gives a finite-field bound.
A run count for the final support covers the remaining sparse regime.
This argument applies to all messages after setup; it is stronger than
a generic-rank assertion alone.

There is also a complete two-permutation, width-two result with
$q=2048$ and the smaller local $[8,4,5]$ source: minimum distance
at least 734004 ($>35\%$), with failure below $2^{-40.22}$.
It supplies a useful intermediate point, but changing the local source
prevents a direct comparison that isolates width.
The remaining grid will hold the source fixed and compare $L=2$ and
$L=3$ separately. An extra final scan will always be identified.
Historical dense-support searches alone do not supply complete grid
entries or upper bounds on the encoders' distances.

\subsection{A coefficient-independent obstruction}\label{sec:group-ceiling}

Coefficient size cannot overcome every limitation of a narrow scan.
The following count gives an upper bound for each fixed encoder.

\begin{lemma}\label{lem:check-group-ceiling}
Let $n=2rm$, and let $H$ be a parity-check matrix with $m$ groups of
$r$ rows each. Suppose each column touches at most $D$ groups, where
$2\le D\le m$. If $D\le b\le m$ and
\begin{equation}\label{eq:group-ceiling-condition}
 2\binom{b-1}{D-1}>\binom{m-1}{D-1},
\end{equation}
then $d(\ker H)\le rb+1$.
\end{lemma}

\begin{proof}
Fix $H$ and sample a uniform set of $b$ row groups. A column touching
$d\le D$ groups is supported entirely in the selected rows with
probability $\binom bd/\binom md\ge\binom bD/\binom mD$.
The expected number of such columns is at least
$n\binom bD/\binom mD$. Condition~\eqref{eq:group-ceiling-condition}
makes this number exceed $rb$. Some selection therefore contains
$rb+1$ columns supported on at most $rb$ rows. These columns are
dependent, and their dependence gives a nonzero word of weight at most
$rb+1$ in $\ker H$.
\end{proof}

The local source has a block diagonal parity check $H_0$, with eight
rows per block. Write $D_j=A_{w,a_j}^{-1}$. With the convention
in~\eqref{eq:width-family}, an output parity check is
\[
 H=H_0P_1^{-1}D_1\cdots P_L^{-1}D_L.
\]
Each $D_j$ has bandwidth $w$ and at most $w+1$ nonzeros per column.
Thus every column of $H$ touches at most $(w+1)^L$ source groups.
Consecutive inverse scans without a permutation have additive
bandwidth: a sequence of widths $u_1,\ldots,u_j$ has at most
$1+\sum_i u_i$ nonzeros per column. This yields tighter counts than
multiplying their individual column weights.

For width one, let $(t_1,\ldots,t_L)$ give the number of consecutive
accumulators after each permutation. We may take
$D=\prod_j(t_j+1)$. Applying the exact integer condition with
$m=131072$ gives
\begin{center}
\begin{tabular}{lrrr}
\toprule
Consecutive scan counts & $D$ & Distance at most & Relative ceiling\\
\midrule
$(1,1)$ & 4 & 832265 & $39.6855\%$\\
$(1,2)$ & 6 & 912849 & $43.5281\%$\\
$(1,1,1)$ & 8 & 949729 & $45.2867\%$\\
$(2,2)$ & 9 & 961553 & $45.8505\%$\\
\bottomrule
\end{tabular}
\end{center}
The percentages are rounded upward. These bounds hold for every
permutation and every coefficient choice, including full-field
coefficients. In particular, two width-one scans cannot reach $44\%$,
even after one additional accumulator without another permutation.
For fixed $D$ and growing $m$, the relative ceiling tends to
$\tfrac12 2^{-1/(D-1)}$. The bound does not determine typical distance:
the best achieved width-one lower bound and its ceiling still leave
a gap between $30.5\%$ and $39.6855\%$.
